\documentclass[11pt]{article}
\usepackage[margin=1in]{geometry}
\usepackage{amsmath,amssymb}
\usepackage{booktabs}
\usepackage{graphicx}
\usepackage[hidelinks]{hyperref}
\usepackage{natbib}
\usepackage{caption}
\captionsetup{font=small,labelfont=bf}
\newcommand{\Wj}{W_{j}}
\newcommand{\Sw}{S_{w,j}}
\newcommand{\Ti}{T_{ij}}

\title{\textbf{Construction of the Cumulative Wetland Effectiveness Variable $W$}\\[3pt]
\large Methods and Literature Review --- from mechanism to a computable, calibrated metric}
\author{MMSD Wetland--Flood Effectiveness Study \quad(\texttt{wetland\_v2})}
\date{July 2026}

\begin{document}
\maketitle
\vspace{-1.1em}
\begin{abstract}\noindent\small
This note documents, end to end, how we construct the cumulative ``wetland effectiveness''
predictor $\Wj$ used in the event--based flood--response models, and---equally---\emph{why} each
modeling choice is made, with every choice anchored to the hydrologic literature. Starting from
the physical mechanisms by which wetlands modify floods (storage, attenuation, delay,
infiltration, connectivity), we adopt the advisor's cumulative form
$\Wj=\sum_i A_i\,S_i\,K(\Ti)$, drop the separate connectivity term as redundant with the
travel--time kernel, and build each remaining factor on a published method: a velocity/celerity
time--of--concentration field for $\Ti$, a Volume--Area storage law and a Cowardin/HGM type
weight for $S_i$, and an exponential travel--time kernel for $K$. We then calibrate the
travel--time engine against observed event lags. The resulting variable is physically defensible
and, reassuringly, reproduces the same peak--shaving signal as the simpler original metric.
\normalsize\end{abstract}

\section{Motivation: turning many wetlands into one number}
The scientific question is whether the configuration of wetlands upstream of a gauge is
associated with the flood response measured at that gauge. This requires collapsing thousands of
heterogeneous wetland polygons into a single, physically meaningful predictor per gauge $j$. The
guiding principle \citep{workflow} is to \emph{avoid inventing variables}: every quantity is
borrowed from an authoritative published method. Following the discussion note, we adopt the
cumulative ``wetland effectiveness'' form
\begin{equation}
\Wj=\sum_{i\in\text{basin}(j)} A_i\,C_i\,S_i\,K(\Ti),
\label{eq:advisor}
\end{equation}
where $A_i$ is the area of wetland $i$, $C_i$ a connectivity/near--stream indicator, $S_i$ a
storage/soil/type modifier, $\Ti$ the travel time from wetland $i$ to gauge $j$, and $K(\cdot)$ a
travel--time decay function. Equation~\eqref{eq:advisor} is a weighted sum: each wetland
contributes its area, up--weighted or down--weighted by how much it can plausibly attenuate the
flood and by how strongly it is connected to the gauge.

\section{Literature review: mechanisms $\rightarrow$ mathematical representation}
Wetlands do not affect floods through a single process, and the literature is explicit that the
\emph{type} and \emph{location} of a wetland govern which process dominates
\citep{acreman2013,bullock2003}. Table~\ref{tab:mech} maps each mechanism to its hydrologic
effect and to the mathematical object we use to represent it; the right--hand column is the
design logic for Eq.~\eqref{eq:advisor}.

\begin{table}[htbp]\centering
\caption{Wetland flood mechanisms and their representation in $\Wj$.}
\label{tab:mech}
\small
\begin{tabular}{p{2.3cm}p{4.3cm}p{6.2cm}}
\toprule
Mechanism & Hydrologic effect & Representation (and literature) \\
\midrule
Storage & Lower runoff volume; reduced peak & storage capacity $V_i$ via Volume--Area
\citep{gleason2007,laneDamico2010,wu2016} \\
Attenuation & Lower peak, flatter hydrograph & travel--time weighting $K(\Ti)$ \citep{acreman2013} \\
Delay & Longer time to peak & travel--time weighting / lag \citep{mcmillan2020} \\
Infiltration & Lower runoff coefficient & soil / wetland--type modifier in $S_i$ \\
Connectivity & Stronger effect for near/connected wetlands & \emph{is} the travel--time weighting
\citep{lane2018,ameli2019} \\
Peak--shaving & Disproportionate reduction in large storms & interaction $W\!\cdot\!P$ (model, not $W$) \\
\bottomrule
\end{tabular}
\end{table}

Three results from this literature drive our specific choices. First, \citet{acreman2013} review
shows attenuation and delay dominate the flood signal (wetlands change \emph{timing and shape}
more than total volume), justifying a travel--time--weighted metric rather than a plain area.
Second, \citet{ameli2019} demonstrate with a physically based model that wetlands within
$\sim$100~m of the channel attenuate peaks far more than wetlands $>$4~km away---i.e. the effect
\emph{decays with hydrologic distance to the network}, which is exactly a travel--time kernel.
Third, \citet{fossey2016} and \citet{evenson2015} show the location effect interacts with type
(riparian wetlands matter most on the main stem; isolated wetlands act as headwater
``gatekeepers'' \citep{golden2016}), motivating a wetland--type modifier $S_i$.

\section{Design decision: dropping the separate connectivity term $C_i$}
Equation~\eqref{eq:advisor} contains both a connectivity indicator $C_i$ (e.g. a $<$100~m
near--stream flag) and a travel--time kernel $K(\Ti)$. These encode the \emph{same} mechanism
twice: a near--stream wetland necessarily has a small $\Ti$ and therefore a large $K(\Ti)$. The
authoritative connectivity framework itself defines connectivity operationally as travel time to
the network \citep{lane2018}, and \citet{ameli2019} express the ``$<$100~m matters most'' finding
as a continuous decay with distance, not a binary. Retaining both double--counts connectivity and
inflates the weight of near--channel wetlands. We therefore drop $C_i$ and let $K(\Ti)$ carry all
of the connectivity information:
\begin{equation}
\boxed{\;\Wj=\sum_i A_i\,S_i\,K(\Ti)\;}
\label{eq:W}
\end{equation}
This is more parsimonious and removes a redundant, hard--to--calibrate parameter.

\section{Constructing each factor}

\subsection{Wetland area $A_i$}
Wetland polygons are the USFWS National Wetlands Inventory (NWI), carrying the Cowardin
classification code. We query them by catchment, deduplicate, and clip; $A_i$ is the clipped
polygon area in an equal--area projection. As a validation, the NWI wetland area of the main
Milwaukee basin (324~km$^2$) matches an independent NLCD estimate (325~km$^2$).

\subsection{Travel time $\Ti$: a velocity/celerity time--of--concentration field}
\paragraph{Why travel time, not distance.} Straight--line distance misranks wetlands because
runoff follows the channel network; a small wetland on the stream can matter more than a large one
off to the side. Travel time along the flow path is the physically correct connectivity measure
\citep{lane2018}. We compute it with the NRCS \emph{velocity method} of time of concentration
\citep{tr55,mccuen}: route the 10~m DEM by D8 flow directions, assign each cell a flow speed, and
sum $L/\text{speed}$ along the path from every cell to the gauge (\texttt{pysheds}
\texttt{distance\_to\_outlet}).

\paragraph{DEM conditioning.} On the low--relief glacial terrain of southeastern Wisconsin, plain
depression filling fragments the drainage. We use WhiteBox least--cost breaching
(\texttt{BreachDepressionsLeastCost}) so the contributing area to each gauge matches its mapped
basin, then \texttt{resolve\_flats} to define D8 directions on flats.

\paragraph{Speed 1 --- celerity, not water velocity.} The original model used the Manning water
velocity $v=\tfrac{1}{n}R^{2/3}S^{1/2}$. But a flood \emph{peak} propagates at the kinematic wave
\emph{celerity}, $c=\beta v$, with $\beta=5/3$ for a wide Manning channel
\citep{lighthill1955,beven2012}. Because we weight wetlands by how their \emph{signal} reaches the
gauge, celerity is the correct speed; using the water velocity systematically overestimates the
travel time. Adopting $c=\tfrac{5}{3}v$ in channel cells de--inflates the modeled travel times by
$\approx$1.5$\times$ (e.g.\ basin--median $479\!\to\!297$~h at the main gauge) while preserving the
spatial pattern.

\paragraph{Speed 2 --- wetland roughness.} With a uniform Manning $n$, a wetland does not actually
slow water in the simulation, so the metric would encode only wetland \emph{position}. Yet
emergent marsh has a much higher roughness ($n\approx0.05$--$0.15$) than upland surfaces
\citep{kadlec1990,arcement1989}. We therefore assign a land--cover--dependent overland resistance,
slowing flow through wetland cells. Now the flow--path travel time \emph{integrates the roughness
of the wetlands the water passes through}: $\Ti$ reflects wetlands actively delaying the flood,
turning $W$ from a position index into a mechanism.

\subsection{Kernel $K(\Ti)$ and the decay scale $\tau$}
We use the exponential kernel $K(T)=e^{-T/\tau}$ (with the harmonic $1/(1+T/\tau)$ as a
robustness variant): it is smooth, bounded, and reduces to the plain wetland fraction when the
weight is constant \citep{workflow}. The decay scale $\tau$ is the travel--time distance over
which a wetland's influence persists. Because absolute travel times on this flat terrain are still
long (hundreds of hours), a fixed $\tau$ (e.g.\ 24~h) would send nearly every weight to zero; we
set $\tau$ adaptively to each catchment's median travel time, which keeps the weighting
well--scaled and comparable across basins, and calibrate it against data (\S\ref{sec:calib}).

\subsection{Storage/type modifier $S_i$}
The advisor's $S_i$ bundles storage, soil, and wetland type. We construct it so that the product
$A_i S_i$ is mechanistically interpretable:
\begin{equation}
S_i = M_{\text{type},i}\times d_i,\qquad A_i S_i = M_{\text{type},i}\,V_i .
\end{equation}
\paragraph{Storage depth $d_i$ (Volume--Area).} Rather than hand--pick coefficients, we use the
published Volume--Area power law $V_i = c\,A_i^{\beta}$
\citep{gleason2007,laneDamico2010,wu2016}, so the depth is $d_i=V_i/A_i=c A_i^{\beta-1}$. We adopt
a glaciated--terrain fit ($c\!=\!0.01725,\ \beta\!=\!1.30$, cf.\ prairie--pothole depressions;
\citealp{wu2016}) appropriate to Wisconsin's glacial landscape, replacing an earlier placeholder.
Then $A_i S_i = M_{\text{type}}V_i$ and
\begin{equation}
\Wj=\sum_i M_{\text{type},i}\,V_i\,e^{-\Ti/\tau},
\label{eq:Wfinal}
\end{equation}
i.e.\ the connectivity--delivered, type--weighted \emph{wetland storage} reaching gauge $j$. This
unifies the storage capacity into $W$: $\Sw=\sum_i V_i$ is the same quantity without the kernel.
\paragraph{Type modifier $M_{\text{type}}$.} Following \citet{acreman2013}, floodplain and emergent
wetlands attenuate peaks more than open--water or lacustrine systems. We map the Cowardin class to
an attenuation weight (emergent $1.0$; forested/scrub--shrub $0.85$; open water/aquatic bed $0.5$;
riverine $0.6$; lacustrine $0.3$)---a transparent, tunable table grounded in the type ordering of
\citet{acreman2013,bullock2003}. Soil (saturated conductivity) is retained as a basin--level
control rather than a per--wetland factor, because point SSURGO values on small polygons are noisy
and infiltration acts at the catchment scale.

\section{Calibrating the travel--time engine}\label{sec:calib}
A model of travel time should be checked against observed response times. From the event pipeline
we have, per gauge, the observed hydrograph lag (time to peak). The SCS relation gives basin
lag $\approx0.6\,T_c$ \citep{tr55}. Regressing observed lag on the modeled travel time across the
nine gauges yields $\text{lag}\approx0.19\,T_c$ ($R^2=0.46$): the model explains about half the
cross--gauge variance, but the median travel time over--represents slow overland/wetland cells
relative to the channel--dominated response, and the \emph{flashy urban creeks fit worst}---their
2--3~h response is driven by imperviousness and storm sewers, not natural travel time. We therefore
report the calibration as a validation and retain the adaptive $\tau$; a fully calibrated,
urbanization--corrected celerity is flagged for future work.

\section{Result and robustness}
The improved variable of Eq.~\eqref{eq:Wfinal}---physical (celerity) travel times, wetland
roughness, the redundant $C$ removed, and a type--weighted glaciated storage $S_i$---reorders the
wetland ranking (the most wetland--dense rural basin now scores highest per unit area). Crucially,
when substituted into the event models it \emph{reproduces the same peak--shaving signal} as the
original connectivity--weighted fraction: the standardized $W\!\cdot\!P$ coefficient on log peak
discharge is $-0.155$ (wild--cluster--bootstrap $p=0.029$) versus $-0.157$ ($p=0.004$) for the
original metric. That the finding is invariant to a substantial, literature--grounded rebuild of
$W$ is itself evidence that it is not an artifact of one arbitrary construction.

\section{What the literature says we still cannot do}
Two honest limits remain. (i) \emph{Identification.} $\Wj$ is gauge--constant and, across only
nine nested catchments, collinear with urbanization; the wetland \emph{level} effect is therefore
not identifiable regardless of how well $W$ is built---only the event--level interaction
$W\!\cdot\!P$ is. A causal wetland level effect needs a regional sample of independent catchments
spanning a wetland gradient decoupled from urbanization. (ii) \emph{Type$\times$location tension.}
\citet{ameli2019} (near--channel wetlands attenuate peaks most) and \citet{fossey2016} (isolated
wetlands most effective at headwaters) are not fully reconciled; a single kernel cannot capture
both, and separating riparian from isolated $W$ components is the natural next refinement.

\begin{thebibliography}{99}\small
\bibitem[Acreman \& Holden(2013)]{acreman2013} Acreman, M., \& Holden, J. (2013). How wetlands
affect floods. \emph{Wetlands}, 33, 773--786.
\bibitem[Ameli \& Creed(2019)]{ameli2019} Ameli, A.~A., \& Creed, I.~F. (2019). Does wetland
location matter when managing wetlands for watershed--scale flood and drought resilience?
\emph{JAWRA}, 55(3), 529--542.
\bibitem[Arcement \& Schneider(1989)]{arcement1989} Arcement, G.~J., \& Schneider, V.~R. (1989).
Guide for selecting Manning's roughness coefficients for natural channels and flood plains.
\emph{USGS Water--Supply Paper 2339}.
\bibitem[Beven(2012)]{beven2012} Beven, K. (2012). \emph{Rainfall--Runoff Modelling: The Primer},
2nd ed. Wiley.
\bibitem[Bullock \& Acreman(2003)]{bullock2003} Bullock, A., \& Acreman, M. (2003). The role of
wetlands in the hydrological cycle. \emph{Hydrology and Earth System Sciences}, 7(3), 358--389.
\bibitem[Evenson et al.(2015)]{evenson2015} Evenson, G.~R., Golden, H.~E., Lane, C.~R., \&
D'Amico, E. (2015). Geographically isolated wetlands and watershed hydrology. \emph{Journal of
Hydrology}, 529, 240--256.
\bibitem[Fossey \& Rousseau(2016)]{fossey2016} Fossey, M., \& Rousseau, A.~N. (2016). Assessing
the long--term hydrological services provided by wetlands under changing climate conditions.
\emph{Journal of Hydrology}, 541, 1287--1302.
\bibitem[Gleason et al.(2007)]{gleason2007} Gleason, R.~A., et al. (2007). Estimating water storage
capacity of existing and potentially restorable wetland depressions. \emph{USGS Open--File
Report 2007--1159}.
\bibitem[Golden et al.(2016)]{golden2016} Golden, H.~E., et al. (2016). Non--floodplain wetlands
affect watershed hydrology: ``gatekeepers'' of downstream flows. \emph{Ecohydrology}.
\bibitem[Kadlec(1990)]{kadlec1990} Kadlec, R.~H. (1990). Overland flow in wetlands: vegetation
resistance. \emph{Journal of Hydraulic Engineering}, 116(5), 691--706.
\bibitem[Lane \& D'Amico(2010)]{laneDamico2010} Lane, C.~R., \& D'Amico, E. (2010). Calculating
the ecosystem service of water storage in isolated wetlands using LiDAR in North Central Florida.
\emph{Wetlands}, 30, 967--977.
\bibitem[Lane et al.(2018)]{lane2018} Lane, C.~R., et al. (2018). Hydrological connectivity of
non--floodplain wetlands to downstream waters: a classification approach. \emph{JAWRA} / US EPA.
\bibitem[Lighthill \& Whitham(1955)]{lighthill1955} Lighthill, M.~J., \& Whitham, G.~B. (1955).
On kinematic waves. I. Flood movement in long rivers. \emph{Proc. R. Soc. Lond. A}, 229, 281--316.
\bibitem[McMillan(2020)]{mcmillan2020} McMillan, H. (2020). Linking hydrologic signatures to
hydrologic processes: a review. \emph{Hydrological Processes}, 34(6), 1393--1409.
\bibitem[McCuen(2009)]{mccuen} McCuen, R.~H. (2009). Uncertainty analyses of watershed time
parameters. \emph{Journal of Hydrologic Engineering}, 14(5), 490--498.
\bibitem[NRCS(1986)]{tr55} Natural Resources Conservation Service (1986). \emph{Urban Hydrology
for Small Watersheds}, Technical Release 55 (TR--55). USDA.
\bibitem[Wu \& Lane(2016)]{wu2016} Wu, Q., \& Lane, C.~R. (2016). Delineation and quantification of
wetland depressions in the Prairie Pothole Region; Volume--Area--Depth storage relations.
\emph{Wetlands}, 36, 215--227.
\bibitem[Discussion note \& workflow(2026)]{workflow} MMSD Wetland--Flood study (2026). Discussion
note and data/modeling workflow (internal).
\end{thebibliography}
\emph{\footnotesize Bibliographic details are representative; verify volume/issue/page before
formal citation.}
\end{document}
